\section{Experimental Setup}

\subsection{Data}

We construct simultaneous speech translation training trajectories from
GigaSpeech \citep{chen2021gigaspeech}, a large-scale multi-domain English
speech corpus. Each example provides English speech and transcript input; we
use the transcript stream as the source-side incremental signal for data
synthesis, and pair it with target translations generated for downstream
SimulST training.

Our main experiments focus on English-to-Chinese (EN$\rightarrow$ZH), where
word-order divergence and future-dependent lexical choices make premature
commitments especially harmful. We additionally evaluate
English-to-German (EN$\rightarrow$DE) and English-to-Japanese
(EN$\rightarrow$JA) to test whether the same synthesis procedure transfers
across typologically different target languages. For all directions, source
utterances are segmented into streaming chunks to simulate incremental ASR
input. The synthesis algorithm then produces a sequence of READ/WRITE
decisions and target deltas for each chunk, as described in
Section~\ref{sec:method}.

\subsection{Training Data Synthesis}

For each source prefix, we sample multiple plausible English continuations
from a base language model. Each sampled future is translated with an
instruction-tuned translation model conditioned on the already committed
target prefix. We then apply token-level distributional consensus over the
next-token candidate sets across futures, committing only tokens supported by
all sampled continuations. The resulting prefix-aligned trajectories are used
to fine-tune the SimulST model, while examples with low quality-estimation
scores are filtered with MetricX-QE before training.

\subsection{SimulST Model and Training}
\label{sec:setup-training}

\paragraph{Backbone.} All SimulST models are obtained by fine-tuning
\textbf{Qwen3-Omni-30B-A3B-Instruct} \citep{qwen3omni}, a Mixture-of-Experts
omni-modal model with $\sim$30B total parameters and $\sim$3B active per
token. We use the same backbone across all data-synthesis variants and all
three language pairs so that observed quality differences are attributable to
the synthesised supervision rather than to model capacity.

\paragraph{Fine-tuning recipe.} We perform parameter-efficient
fine-tuning with LoRA \citep{hu2022lora} of rank $r{=}32$ and scaling
$\alpha{=}32$ applied to all linear projections; the audio encoder (ViT) and
the modality aligner are frozen, while the LLM expert layers receive LoRA
updates. Training is run with the Megatron-Swift toolkit using packing,
expert parallelism of size $4$, fused MoE permutation, and grouped GEMM with
auxiliary load-balancing loss (coefficient $10^{-3}$). We optimize with
AdamW (clip $1.0$, weight decay $0.01$) at peak learning rate $10^{-4}$ with
linear warmup over $5\%$ of total steps and a cosine decay to $10^{-5}$.
Each example is packed to a maximum sequence length of $2048$ tokens, with
micro-batch size $1$ and global batch size $4$; we train for one epoch.

\paragraph{Audio chunking.} For training, source audio is segmented into
$960$\,ms chunks at $16$\,kHz; each utterance is assigned a uniformly-random
chunk \emph{multiplier} $r$ which groups $r$ consecutive $960$\,ms chunks
into a single streaming step. The corresponding aligned target tokens are
grouped accordingly. This randomised multiplier teaches the model to
operate at a range of streaming granularities without requiring separate
checkpoints per latency target, mirroring the data layout used by
EAST~\citep{EAST}.

\paragraph{Training data scale.} For each synthesis policy we randomly
downsample the QE-filtered pool to $12{,}500$ trajectories so that
comparisons across policies are matched on training-set size; the original
pool sizes (before downsampling) are reported in
Section~\ref{sec:results}.

\paragraph{Compute.} Each fine-tuning run uses $4{\times}$ NVIDIA L40S
($48$\,GB) GPUs and completes in roughly one wall-clock hour for the
EN$\rightarrow$ZH variants and up to seven hours for the EAST training data
on EN$\rightarrow$DE / EN$\rightarrow$JA, where the supervision sequences
are longer.

\subsection{Inference and Evaluation Protocol}
\label{sec:setup-eval}

\paragraph{Streaming inference.} At inference time we wrap the fine-tuned
SimulST model in a custom \texttt{simuleval}~\citep{ma2020simuleval} agent
backed by vLLM~\citep{kwon2023vllm} (tensor-parallel size $2$, max model
length $32{,}768$, eager execution). Decoding uses temperature $\tau{=}0.6$,
top-$p{=}0.95$, top-$k{=}20$, and at most $30$ new tokens per streaming
step; the agent enforces a minimum start latency of $2$\,s to mimic the
end-of-segment delay of a streaming ASR front-end. To measure the
quality--latency trade-off of each model with a single set of weights, we
sweep the source segment size over
$\{960, 1920, 2880, 3840\}$\,ms, holding all other decoding hyper-parameters
fixed.

\paragraph{Test set.} We evaluate on the ACL~6060 development
set~\citep{salesky2023acl6060}, which provides multi-way parallel
human-translated references for EN$\rightarrow$\{ZH, DE, JA\} aligned to
long-form English talk audio. This setting is substantially harder than
sentence-level test sets because each utterance is a multi-minute talk
segment, exposing the model to the full streaming regime rather than to
isolated sentences.

\paragraph{Quality metrics.} For each segment-size point we report two
quality numbers. \emph{Chunk-level BLEU} is computed by SimulEval over the
streaming hypothesis using sacre\-BLEU~\citep{post2018sacrebleu} with
language-specific tokenisers
(\texttt{zh}, \texttt{13a}, \texttt{ja-mecab} for ZH, DE, JA respectively).
\emph{Long-form BLEU and COMET} are computed by aligning the streaming
hypothesis to the reference sentence boundaries with
\texttt{mwerSegmenter}~\citep{matusov2005evaluating} and then scoring with
sacre\-BLEU and \textsc{XCOMET-XL}~\citep{guerreiro2024xcomet}.
The long-form numbers are our primary quality measure because they correctly
penalise hypotheses that drift across sentence boundaries.

\paragraph{Latency metric.} Following recent SimulST
work~\citep{papi2023longyaal}, we report \emph{LongYAAL}, a
length-adaptive average lagging variant designed for long-form streaming,
in its \emph{computation-unaware} (CU) form so that latency reflects the
model's policy decisions rather than transient inference-engine cost.
Reported latency is the mean LongYAAL\,(CU) in milliseconds across the
ACL~6060 dev talks at each segment size.

\paragraph{Baselines.} We compare against fixed-latency policies
(wait-$k$~\citep{ma2019stacl}), alignment-based prefix construction
(EAST~\citep{EAST}), local-agreement decoding
(LA-$N$~\citep{liu2020low}), and recent LLM-based data-synthesis baselines
(Simul-MuST-C, TAF, InfiniSST). All baselines are fine-tuned from the same
Qwen3-Omni-30B-A3B-Instruct backbone with the recipe described in
Section~\ref{sec:setup-training}, so any quality difference reflects the
synthesised supervision rather than backbone or optimiser changes. Our
primary comparison is on EN$\rightarrow$ZH; EN$\rightarrow$DE and
EN$\rightarrow$JA serve as cross-lingual transfer settings.
